\documentclass[10pt,a4paper]{article}
\usepackage[utf8]{inputenc}
\usepackage[a4paper,left=3cm,right=2cm]{geometry}
\usepackage{amsmath}
\usepackage[thmmarks, amsmath, thref]{ntheorem}
\usepackage{amsfonts}
\usepackage{amssymb}
\usepackage{fancyhdr}
\pagestyle{fancy}
\usepackage{anysize}
\usepackage{graphicx}
\usepackage{float}
\usepackage{hyperref}
\usepackage{pgfplots}
\pgfplotsset{compat=1.17}
\usepackage{enumitem}
\newtheorem{theorem}{Theorem}
\theoremsymbol{\ensuremath{\blacksquare}}
\newtheorem*{solution}{Solución:}

\usepackage{xcolor}

% Margenes y formalidades

\textwidth 6.25in % Margen Derecho y ancho del texto
\textheight 8.5in
\oddsidemargin 0in % Margen Izquierdo
\headheight 0in % Largo del texto
\leftmargin 1in
\parindent 0em % Salto entre la indentación (Sangría)https://www.overleaf.com/4519127669sdpptkhvsbfn
\topmargin -0.5in % Margen Superior
\parskip 0.5ex % Salto entre parrafos
\baselineskip 1pt



\lhead{
\includegraphics[scale=0.11]{figuras/logoUSM.png}} % texto izquierda de la cabecera
%\chead{}                                      % texto centro de la cabecera
\rhead{Universidad Técnica Federico Santa María\\Ingeniería Comercial } % número de página a la derecha
\renewcommand{\headrulewidth}{0.4pt} % grosor de la línea de la cabecera
\renewcommand{\footrulewidth}{0.4pt} % grosor de la línea del pie

\begin{document}
\vspace{5cm}

\begin{center}
\begin{Large}
\textbf{\\Formulario Econometría ICS-294}
\end{Large}\\
\end{center}

    
\section{Formulario}
\begin{itemize}

\item El intervalo de confianza de $100(1-\alpha)\%$ para la media poblacional $\mu$ cuando $\sigma$ es \emph{\textbf{conocido}}:

$$\color{red}{\bar{x} \pm z_{\frac{\alpha}{2}}\frac{\sigma}{\sqrt{n}}}$$

\item El intervalo de confianza de $100(1-\alpha)\%$ para la diferencia de medias
$\mu_x-\mu_y$ es:

$${\color{red}(\bar{X}-\bar{Y}) \pm z_{\frac{\alpha}{2}}\sqrt{\frac{\sigma_x^2}{n_x}+\frac{\sigma_y^2}{n_y}}}$$


\item El intervalo de confianza de $100(1-\alpha)\%$ para la media poblacional $\mu$ cuando $\sigma$ es \emph{\textbf{desconocido}}:

$$\color{red}{\bar{x} \pm z_{(n-1)\frac{\alpha}{2}}\frac{S}{\sqrt{n}}}$$

\item El intervalo de confianza de $100(1-\alpha)\%$ para la proporción poblacional $p$ es:

$$\color{red}{\hat{p} \pm z_{(n-1)\frac{\alpha}{2}}\sqrt{\frac{\hat{p}(1-\hat{p})}{{n}}}}$$


\item  Estimadores MCO para RLS.
{\color{red}
$$\hat{\beta}_0 = \bar{y} - \hat{\beta}_1 \bar{x} $$
$$\hat{\beta}_1 = \frac{\sum_{i=1}^{n}(x_i - \bar{x})(y_i-\bar{y})}{\sum_{i=1}^{n}(x_i - \bar{x})^2}=\frac{Cov(x,y)}{Var(x)}$$}


\item  Estimadores MCO para RLM.
{\color{red}
$${\hat{\beta}} = (X^TX)^{-1} X^TY $$}

    
\item SST la suma de cuadrados totales (total sum of squares):
	{\color{red}	\[ SST = \sum_{i=1}^{n} (y_i - \bar{y})^2 \]}
		
		\item  SSE la suma de cuadrados explicados (explained sum of squares):
	{\color{red}	\[ SSE = \sum_{i=1}^{n} (\hat{y}_i - \bar{y})^2 \]}
		
		\item  SSR la suma de cuadrados no explicados (residual sum of squares):
	{\color{red}	\[ SSR = \sum_{i=1}^{n} \hat{u}_i^2 \]}

		\item  $R^2$:
	{\color{red}	\[ R^2= 1 - \frac{SSR}{SST} \]}

    		\item  $R^2$ ajustado:
	{\color{red}	\[ R^2_{adj}= 1 - \frac{n-1}{n-k-1}\frac{SSR}{SST} \]}


\end{itemize}

\newpage

\begin{figure}
    \centering
    \includegraphics[width=1.2\linewidth, angle=90]{figuras/normal_critic.png}

    
\end{figure}

\end{document}

   